\documentclass[11pt,a4paper]{article}
\usepackage[margin=1in]{geometry}
\usepackage{amsmath,amssymb,amsthm,braket,booktabs,longtable,hyperref}
\usepackage[T1]{fontenc}
\usepackage{lmodern}
\title{Q-SENTINEL D1: Mathematical Model and Protocol Specification}
\author{Quantum / Detection Workstream}
\date{2026}

\newtheorem{theorem}{Theorem}
\newtheorem{lemma}{Lemma}
\newtheorem{definition}{Definition}
\newtheorem{corollary}{Corollary}

\begin{document}
\maketitle

\begin{abstract}
This document is the mathematical specification for the D1--D6 Q-SENTINEL
prototype. It fixes the six-state Pauli alphabet, Bell-pair resource,
teleportation measurement/correction map, verification predicate, basis-guessing
forgery bound, Hoeffding/Chernoff calibration quantities, and the observables
used by the six deterministic detectors. The probability statements are
proved under the explicit independent-round and basis-guessing model. A real
security proof for a complete deployed QDS protocol would additionally require
an authenticated classical channel, a composable security definition, and
independent peer review.
\end{abstract}

\section{Notation and State Alphabet}
Let $\mathcal{B}=\{X,Y,Z\}$ denote the Pauli measurement bases. The six allowed
public-key states are
\begin{align*}
Z &: \ket{0},\ket{1},\\
X &: \ket{+}=\frac{\ket{0}+\ket{1}}{\sqrt2},\quad
     \ket{-}=\frac{\ket{0}-\ket{1}}{\sqrt2},\\
Y &: \ket{+i}=\frac{\ket{0}+i\ket{1}}{\sqrt2},\quad
     \ket{-i}=\frac{\ket{0}-i\ket{1}}{\sqrt2}.
\end{align*}
For each round $j\in\{1,\ldots,n\}$, a basis $B_j\in\mathcal B$ and an
associated eigenstate label $E_j$ are generated from a classical seed. The
expected projective outcome is denoted $e_j\in\{0,1\}$, with plus/0 eigenstates
mapped to $0$ and minus/1 eigenstates mapped to $1$.

\section{Bell Resource and Teleportation}
The shared resource is
\[
\ket{\Phi^+}_{BC}=\frac{\ket{00}+\ket{11}}{\sqrt2}.
\]
Let the input qubit be
\[
\ket{\psi}_A=\alpha\ket0+\beta\ket1,\qquad
|\alpha|^2+|\beta|^2=1.
\]
The standard Bell-state measurement (BSM) on $A,B$ yields two classical bits
$m_1m_2$. The receiver's state before correction is, up to a global phase,
\begin{equation}
\ket{\psi_C^{\rm raw}}=P_{m_1m_2}\ket\psi,
\end{equation}
where
\[
\begin{array}{c|c}
 m_1m_2 & P_{m_1m_2} \\ \hline
 00 & I\\
 01 & X\\
 10 & Z\\
 11 & XZ
\end{array}
\]
and the receiver applies the same Pauli operator to undo the BSM-dependent
rotation. Therefore
\begin{equation}
P_{m_1m_2}^\dagger P_{m_1m_2}\ket\psi=\ket\psi.
\end{equation}
Since Pauli operators are unitary, the correction preserves normalization and
fidelity.

\begin{lemma}[Uniform BSM outcomes]
For a maximally entangled $\ket{\Phi^+}$ resource, each Bell-basis outcome has
probability $1/4$, independently of the input pure state.
\end{lemma}
\begin{proof}
The teleportation decomposition is
\[
\ket\psi\ket{\Phi^+}
=\frac12\sum_{k=1}^4 \ket{\mathrm{Bell}_k}\,P_k\ket\psi.
\]
Each branch has amplitude norm $1/2$, hence probability $1/4$.
\end{proof}

\section{QDS Key Material and Public-Key Distribution}
A seed expands deterministically to
\[
K=((B_1,E_1),\ldots,(B_n,E_n)),
\qquad B_j\in\{X,Y,Z\}.
\]
The corresponding eigenstates are teleported through the Bell-pair resource.
The software records pair-provenance identifiers and a commitment
\[
C=H(\mathrm{verifier\_id}\parallel B_1\parallel\cdots\parallel B_n).
\]
A signature envelope contains the message digest $H(M)$, commitment $C$, a
fresh nonce, verifier identity, and pair-provenance identifiers.

\section{Verification Predicate}
For round $j$, the verifier performs a projective measurement in the committed
basis $B_j$ and obtains $M_j\in\{0,1\}$. Define
\[
D=\sum_{j=1}^n \mathbf 1[M_j\neq e_j],
\qquad r=\frac{D}{n}.
\]
In the noiseless honest model each distributed public-key state is exactly an
eigenstate of its verification basis, so
\[
D=0 \quad\text{with probability }1.
\]
The D3 reference implementation therefore accepts an untampered signature if
and only if all message, identity, commitment, provenance, freshness, and
projective-measurement predicates pass the configured mismatch threshold.

\begin{theorem}[Deterministic honest acceptance]
In the ideal noiseless model, with matching message digest, verifier identity,
basis commitment, pair provenance and fresh nonce, the honest signature is
accepted with probability $1$ when the mismatch threshold is zero.
\end{theorem}
\begin{proof}
Each measurement is performed in the eigenbasis of the distributed state, so
$M_j=e_j$ deterministically for every $j$. Hence $D=0$. All classical binding
predicates are equalities by construction, and freshness holds for a new nonce.
Therefore every term in the acceptance predicate is true.
\end{proof}

\section{Basis-Guessing Forgery Bound}
Assume an attacker does not know the committed basis sequence and that rounds
are independent. For a two-basis alphabet $\{X,Z\}$, the attacker chooses the
correct basis with probability $1/2$ and otherwise measures in the wrong basis,
which produces a random eigenvalue. Even in the wrong-basis case the attacker
can guess the expected sign with probability $1/2$, so
\[
p_2=\frac12(1)+\frac12\left(\frac12\right)=\frac34.
\]
For the six-state alphabet $\{X,Y,Z\}$,
\[
p_6=\frac13(1)+\frac23\left(\frac12\right)=\frac23.
\]

\begin{theorem}[Independent-round forgery bound]
Under the basis-guessing model above, the probability that a forger passes all
$n$ rounds satisfies
\[
P_{\mathrm{forge}}(n)\le p^n,
\qquad p\in\left\{\frac34,\frac23\right\}.
\]
\end{theorem}
\begin{proof}
The probability of passing a single round is at most $p$. Independence of the
rounds gives the product bound $p^n$.
\end{proof}

\begin{corollary}[Numerical values]
At $n=128$,
\[
\left(\frac34\right)^{128}\approx 1.0182\times10^{-16},
\]
and
\[
\left(\frac23\right)^{128}\approx 2.8862\times10^{-23}.
\]
At $n=256$ the corresponding values are approximately
$1.0368\times10^{-32}$ and $8.3299\times10^{-46}$.
The first 128-round value is the explicit D3 benchmark quoted in the project
dossier.
\end{corollary}

\section{Hoeffding and Chernoff Statistical Bounds}
Let $X_1,\ldots,X_n\in\{0,1\}$ be mismatch indicators with mean $p$. Hoeffding's
one-sided inequality gives
\[
\Pr\left(\hat p-p\ge\varepsilon\right)\le e^{-2n\varepsilon^2}.
\]
For a target one-sided error budget $\alpha$,
\[
\varepsilon_\alpha=\sqrt{\frac{\ln(1/\alpha)}{2n}},
\qquad
p_{\rm th}=p+\varepsilon_\alpha.
\]
The Bernoulli Chernoff--KL form is
\[
\Pr(\hat p\ge q)\le
\exp\left[-nD_{\mathrm{KL}}(q\Vert p)\right],
\]
with
\[
D_{\mathrm{KL}}(q\Vert p)=q\ln\frac qp+(1-q)\ln\frac{1-q}{1-p}.
\]
The D5 implementation exposes both bounds. For finite-sample validation it
also computes the exact binomial operating probability; the exact tail should
not be confused with the generally looser concentration upper bound.

\section{Channel and Entanglement Observables}
Define
\[
\mathrm{QBER}=\frac{D}{n}.
\]
The Bell-state integrity monitor uses the CHSH value $S$ and Bell-state fidelity
\[
F_{\Phi^+}=\bra{\Phi^+}\rho\ket{\Phi^+}.
\]
The classical CHSH limit is $S\le2$ and the quantum Tsirelson bound is
$S\le2\sqrt2$. The D4 channel detector combines QBER, a four-outcome $\chi^2$
goodness-of-fit test on BSM counts, a one-sided CUSUM, and Wald's SPRT.

\section{Freshness, Identity and Provenance}
Every signature carries a nonce. A consumed nonce cannot be reused. The D4
replay guard maintains an explicit registry. The identity-binding detector
checks verifier identity, basis commitment, message binding and Bell-pair
provenance. These predicates can be mirrored by the team's permissioned ledger
without importing a learned model into the trust path.

\section{Six-Detector Mapping}
\begin{longtable}{lll}
\toprule
Detector & Attack / purpose & Primary evidence\\
\midrule
D1 & Baseline integrity & eigenstate measurement mismatch\\
D2 & Signature forgery & Hoeffding/Chernoff mismatch bound\\
D3 & Entanglement hijack & CHSH and Bell fidelity\\
D4 & Intercept-resend / collective attacks & QBER, $\chi^2$, CUSUM, SPRT\\
D5 & Replay attacks & nonce freshness\\
D6 & Impersonation / unauthorized verification & identity, commitment, provenance\\
\bottomrule
\end{longtable}

\section{Implementation and Security Boundary}
The D2--D6 code is a simulation-first engineering artifact. Its deterministic
rules can be audited and reproduced, but the repository does not claim a
composable end-to-end security proof for all protocol environments. A
publication-grade QDS security proof must additionally specify the authenticated
classical channel, adversary capabilities, key distribution assumptions, finite
sample effects, noise model, multi-verifier composition, and a formal security
notion. Those assumptions are deliberately visible here rather than hidden.

\end{document}
